% Alternate brick builds from a real inventory, with guarantees by construction.
% Draft skeleton (2026-10-08). Every number carries a "% src:" comment naming the results file it comes from,
% so it can be checked; "% src: ... (pending merge)" means the file is on a branch not yet merged into main.
% Tables are generated: run `python paper/tables.py` from the repository root after
% `python -m brickwfc.experiments.paper_tables`.
\documentclass[10pt,twocolumn]{article}

\usepackage[utf8]{inputenc}
\usepackage[T1]{fontenc}
\usepackage{lmodern}
\usepackage[margin=0.8in]{geometry}
\usepackage{amsmath,amssymb}
\usepackage{graphicx}
\usepackage{booktabs}
\usepackage{xcolor}
\usepackage{url}
\usepackage[hidelinks]{hyperref}

\newcommand{\VENUE}{Draft --- venue to be decided}
\newcommand{\todo}[1]{\textcolor{red}{[TODO: #1]}}
% A figure that falls back to a framed placeholder until its image exists.
\newcommand{\figorplaceholder}[3]{%
  \IfFileExists{#1}{\includegraphics[width=#2]{#1}}%
  {\fbox{\parbox[c][7em][c]{\dimexpr#2-2\fboxsep-2\fboxrule\relax}{\centering\small\textit{Placeholder: #3}\\\texttt{\detokenize{#1}}}}}}

\title{Alternate Brick Builds from a Real Inventory,\\ with Stability, Inventory and Shape Guaranteed by Construction}
\author{\todo{authors and affiliation}\\ \small\VENUE}
\date{}

\begin{document}
\maketitle

\begin{abstract}
Given the parts of a real brick set (any set number) and optionally a target shape, we generate new alternate
builds that use only those parts, stand up on a table, form one piece, and resemble official designs. The
physical and inventory requirements are guaranteed by construction rather than filtered afterwards: a bottom-up
decoder places one part at a time under masks for support, inventory counts, the target shape, balance (no
tip-over) and connectivity. Style is learned from 28 official building sets, as contact statistics that weight the
decoder's choices and as a size-matched judge that cannot be won by repetition. A generate, verify, dedupe, judge,
select funnel returns the best candidates. We set out to push these constraints into the propagation step of Wave
Function Collapse (WFC); the constrained WFC is valid by construction too, but a bottom-up decoder matches or beats
it on every measure we took, at 10--100$\times$ lower cost, and WFC fails to finish at the scale of real
inventories. We report this comparison throughout, and the alternate-build system as the result.
\end{abstract}

% ------------------------------------------------------------------------------------------------------------
\section{Introduction}

A brick set comes with one official design, but its parts can build many others. Designing a good alternate build
by hand takes hours; a generator has to respect three hard requirements at once: use only the parts in the box,
stay physically sound (every part supported, nothing tipping over, one liftable piece), and, if the user wants a
particular object, stay inside its shape. On top of these, the result should look like something a designer made.

Text-to-brick generation \cite{pun2025brickgpt} produces stable structures from unlimited bricks of a few sizes;
legolization \cite{luo2015legolization} turns a given 3D model into a stable layout. Neither starts from a fixed,
real inventory, which is the situation of anyone holding a set.

We treat the requirements as constraints on a sequential decoder and enforce each one by construction, so that every
candidate the generator emits is valid. Our contributions:
\begin{enumerate}
  \item \textbf{Three guarantees by construction} for alternate builds from a real inventory: support (every part
    rests on a stud), inventory (never more of a part than the set has) and target shape (nothing outside it),
    plus two global ones added during the work: balance (no stud-connected cluster's centre of mass outside its
    ground contact, checked exactly) and one piece (a repair pass that keeps the build a single connected cluster).
  \item \textbf{A measured comparison of where to enforce them}: inside WFC's propagation, as planned, versus in a
    bottom-up decoder. Both are valid by construction; the bottom-up decoder wins on validity within a time budget,
    inventory use, physical robustness, shape fill and speed, and WFC does not finish on full real inventories
    (Section~\ref{sec:experiments}).
  \item \textbf{Learned style from official sets}: contact statistics from 28 building sets used as a proposal in
    the decoder, and a judge (a divergence of contact statistics) compared at matched size, which ranks repetitive
    builds below every real model.
  \item \textbf{A complete pipeline}: generate, verify, dedupe, judge, select, with step-by-step building
    instructions as output, from any set number (an LDraw model if one exists, otherwise the set's inventory from
    Rebrickable, whose inventories agree with the models at 0.953 weighted Jaccard on 28 sets).
    % src: results/m12c_rebrickable.md (pending merge): per (part, color) mean 0.953
\end{enumerate}

% ------------------------------------------------------------------------------------------------------------
\section{Related work}

\paragraph{Wave Function Collapse and model synthesis.}
WFC \cite{gumin2016wfc} generates tile grids that locally resemble an example by repeatedly collapsing the
lowest-entropy cell and propagating adjacency constraints; it descends from model synthesis \cite{merrell2007model}.
Nested WFC \cite{nie2023nwfc} addresses conflicts and cost for large outputs. Our multi-cell parts are WFC tiles with
integrity rules, and non-pairwise constraints (support, inventory, shape) prune domains during propagation.

\paragraph{Legolization and stability.}
Legolization \cite{luo2015legolization} converts a 3D model into a brick layout and refines weak regions with a
force-based stability analysis. StableLego \cite{liu2024stablelego} analyses block-stacking stability by force
balance and releases layouts with stability labels. Our physics check is a linear program over vertical forces with
push-only contacts and pull-limited stud clutch (Section~\ref{sec:physics}); unlike both, we do not start from a
target model with free choice of bricks, but from a fixed inventory.

\paragraph{Text-to-brick generation.}
BrickGPT \cite{pun2025brickgpt} fine-tunes a language model on StableText2Brick \cite{stabletext2brick}, structures
legolized from ShapeNet \cite{chang2015shapenet} objects in a $20^3$ grid, and enforces stability by rolling back
unstable bricks during decoding.
% src: CLAUDE.md "Positioning vs BrickGPT" -- their reported 98.8% stable, ~40 s/build on GPU: \todo{verify against their paper}
We differ in the input (an inventory and an optional shape instead of text), the source of style (official sets
instead of legolized meshes), and how constraints are met (masks at decode time, so no rollback), and we run on a
CPU without training for the statistics proposal.

% ------------------------------------------------------------------------------------------------------------
\section{Method}

\subsection{Parts and builds}
A build lives on a grid of cells, one stud wide and one plate tall, indexed $(y,z,x)$. The catalog holds 63
rectangular parts: 24 plates, 15 bricks (three plates tall), 9 tiles (no studs) and 15 lookalikes that fill the same
box with the same studs (round plates, grille bricks), kept as their own parts so that a build uses the set's actual
pieces.
% src: results/m8_parts.md ("Catalog ... 63 rectangular parts", 24/15/9/15)
Parts are read from LDraw files \cite{ldraw}: an official model from the LDraw Official Model Repository
\cite{ldrawomr} is voxelized onto the stud grid (parts at 90-degree rotations and grid-aligned positions), and its
catalog parts form the inventory. With the catalog, 51\% of the parts of 20 OMR sets are usable, up from 8\% with
four plate sizes.
% src: results/m8_parts.md ("The usable share goes from 8% to 51%")
\todo{M12b-4 slopes and corner parts, if merged before submission.}

\paragraph{Any set number.} Sets without an LDraw model fall back to Rebrickable's public inventories
\cite{rebrickable}: the latest inventory version without spares, part numbers mapped to catalog parts directly or
through print and mold relationships, and colors mapped to LDraw codes. On the 28 sets of Section~\ref{sec:style}
the two sources agree at a weighted Jaccard index of 0.953 per (part, color), 0.969 per part; the differences are
lookalike numbering, printed parts, and two sets where the sources disagree.
% src: results/m12c_rebrickable.md (pending merge): 0.953 mean (median 0.981), per part 0.969, per color 0.964

\subsection{Bottom-up decoder with masks}
The decoder fills the grid layer by layer from the ground. At each visited cell it draws an order over the candidate
(part, color, orientation, offset) choices from a proposal distribution and places the first that passes every mask:
\begin{itemize}
  \item \textbf{integrity}: inside the grid, overlapping nothing, nothing on a studless tile;
  \item \textbf{support}: above the ground, at least one stud under its bottom layer;
  \item \textbf{inventory}: a (part, color) whose count is used up is skipped;
  \item \textbf{shape}: every cell inside the target mask;
  \item \textbf{balance}: the stud-connected cluster the part joins keeps its centre of mass inside the convex
    hull of its ground contacts (Section~\ref{sec:physics});
  \item \textbf{one piece}: a repair pass keeps the largest cluster, returns the other clusters' parts to the
    inventory and refills above the ground, where every new part rests on the kept cluster.
\end{itemize}
A part may cover the visited cell with any of its cells (``any offset''), which lets it reach into a gap from the
side; this variant (D2) is our default. Every random choice uses a seeded generator, so a seed determines a build.

\paragraph{Constrained WFC (compared method).} The same parts become WFC tiles (one per cell of a part, per
orientation and color) with integrity rules as pairwise adjacencies. Support, inventory and shape are constraint
objects that prune domains after arc consistency, alternating to a fixpoint; inventory needs bounded chronological
backtracking because the count limit is global.

\subsection{Physics}\label{sec:physics}
A build stands on a table if a linear program over vertical forces is feasible: each part has weight proportional to
its cells at its footprint's centre, contacts can push anywhere on shared cells, studs can pull up to a clutch limit,
and every part balances forces and both horizontal torques; the table pushes but never pulls. The result is the
required clutch per stud (solved with HiGHS \cite{huangfu2018highs}). With unbounded clutch, stud-joined parts act as one rigid cluster, so a build needs finite
clutch exactly when every cluster's centre of mass lies within the convex hull of its ground contacts; we check this
in closed form (union-find and a 2D hull), which agrees with the LP on 240 of 240 candidates and lets the decoder
enforce it incrementally.
% src: results/m12_pipeline.md (step 1b: "agrees with the LP on 240/240 pipeline candidates")

\subsection{Proposal: learned style}\label{sec:style}
From 28 official building sets (12,595 parts on the stud grid) we count contacts: which part type rests on which,
which sits beside which, how often touching parts share a color, and how many parts each part rests on.
% src: results/m10_style.md ("28 building sets ... 12,595 parts on the stud grid")
The decoder's proposal weights a candidate by the likelihood ratio of its contacts with the part below and the parts
beside the cell. A learned sequence model can replace or multiply this proposal under the same masks; we train a
small transformer \cite{vaswani2017attention} on brick token sequences from StableText2Brick and the official sets.
\todo{results of the learned proposal (M12b step 2), pending training}
% src: results/m12b_tokens.md (pending merge): 42,604 + 4,785 BrickGPT structures, 168 set sequences, vocab 488

\subsection{Judge and funnel}
The judge is minus the mean Jensen--Shannon divergence \cite{lin1991divergence} between a build's contact statistics
and the learned ones (vertical pairs, side-by-side pairs, support counts), plus the gaps in same-color rates. The
divergence shrinks with size (Spearman $\rho=0.93$ between judge and size over the real models), so builds and real
models are compared after subsampling both to the same number of contact events. A per-contact likelihood would be
gamed by repetition; under the size-matched divergence every repetitive build we tried (towers, sheets, walls) scores
below every real model of its size.
% src: results/m10_style.md ("Spearman rho = 0.93"; degenerate builds table)
The funnel generates $n$ candidates, verifies them (all masks plus physics and a size floor), removes near-duplicates
in a feature space (per-layer fill, part and color shares, a pooled height map), judges the rest, and selects a
diversity-aware top-$k$; each selected build is written as an LDraw file with one building step per layer.

% ------------------------------------------------------------------------------------------------------------
\section{Experiments}\label{sec:experiments}

Generators: \textbf{A} plain WFC with the constraint checked afterwards; \textbf{B} WFC with the constraints in
propagation; \textbf{D} the bottom-up decoder (min corner on the visited cell); \textbf{D2} the decoder with any
offset; and the full pipeline (D2 with learned style, balance and one piece). Table~\ref{tab:main} collects the main
numbers and Table~\ref{tab:ablation} the ablations; each cell names its source experiment, and numbers are comparable
within a column only when they share it.

\begin{table*}[t]
  \centering\small
  \resizebox{\textwidth}{!}{% Generated by paper/tables.py from results/m13_tables.md; do not edit.
\begin{tabular}{lrrrrrrrrr}
\toprule
\textbf{generator} & \rotatebox{70}{valid \%} & \rotatebox{70}{stable \% (rule)} & \rotatebox{70}{stands \% (physics)} & \rotatebox{70}{utilization \%} & \rotatebox{70}{fill / IoU \%} & \rotatebox{70}{s / accepted design} & \rotatebox{70}{distinct designs \%} & \rotatebox{70}{norm. mean distance} & \rotatebox{70}{timeout \%} \\
\midrule
WFC + post-filter (A) & 1.5\textsuperscript{\,m9} & 7.8\textsuperscript{\,m7} & 73.5\textsuperscript{\,m9} & 83.8\textsuperscript{\,m9} & -- & 31.9\textsuperscript{\,m9} & -- & -- & -- \\
WFC + constraints (B) & 2.5\textsuperscript{\,m9b} & 100.0\textsuperscript{\,m7} & 1.2\textsuperscript{\,m9b} & 80.8\textsuperscript{\,m9b} & 73.0\textsuperscript{\,m11} & 1.19$\times 10^{3}$\textsuperscript{\,m9b} & -- & -- & 98\textsuperscript{\,m9b} \\
Bottom-up (D) & 100.0\textsuperscript{\,m9b} & 100.0\textsuperscript{\,m7} & 60.0\textsuperscript{\,m9b} & 99.5\textsuperscript{\,m9b} & 68.0\textsuperscript{\,m11} & 0.319\textsuperscript{\,m9b} & -- & -- & 0\textsuperscript{\,m9b} \\
Bottom-up, any offset (D2) & 100.0\textsuperscript{\,m9b} & 100.0\textsuperscript{\,m9b} & 83.5\textsuperscript{\,m9b} & 99.5\textsuperscript{\,m9b} & 77.7\textsuperscript{\,m11} & 0.992\textsuperscript{\,m9b} & -- & -- & 0\textsuperscript{\,m9b} \\
D2 + learned style (M10) & 100.0\textsuperscript{\,m10} & -- & -- & 99.4\textsuperscript{\,m10} & -- & 3.26\textsuperscript{\,m10} & -- & -- & 0\textsuperscript{\,m10} \\
Pipeline: D2 + style + balance + one piece & 100.0\textsuperscript{\,dm} & 100.0\textsuperscript{\,dm} & 100.0\textsuperscript{\,dm} & 79.5\textsuperscript{\,dm} & 64.0\textsuperscript{\,dm} & 45.3\textsuperscript{\,dm} & 71\textsuperscript{\,dv} & 1.19\textsuperscript{\,dv} & 0\textsuperscript{\,dm} \\
Pipeline with WFC as generator & 0.0\textsuperscript{\,dm} & 0.0\textsuperscript{\,dm} & -- & -- & -- & -- & -- & -- & 100\textsuperscript{\,dm} \\
D2 + learned proposal (M12b) & \textit{pending}\textsuperscript{\,pr} & -- & -- & \textit{pending}\textsuperscript{\,pr} & -- & \textit{pending}\textsuperscript{\,pr} & -- & -- & -- \\
BrickGPT structures (their data) & -- & -- & \textit{pending}\textsuperscript{\,bg} & -- & -- & -- & -- & -- & -- \\
\bottomrule
\end{tabular}
}
  \caption{Main results per generator. Superscripts name the source experiment, and setups differ: m6/m7 plates
    without inventory; m9 plates of 10 sets; m9b the full catalog of 10 sets; m10 28 sets, held-out folds; m11 10 sets
    $\times$ 5 shapes; dm the spaceship demo from 6 sets; bl/cn the balance and one-piece studies; dv the diversity
    pass; pr the learned-proposal comparison; bg BrickGPT's structures. Percent columns are shares of attempted builds; utilization and fill are over finished
    builds. Diversity: effective share of distinct designs among verified candidates, distances in units of the
    setup's typical D2 distance.}
  \label{tab:main}
\end{table*}

\begin{table*}[t]
  \centering\small
  \resizebox{\textwidth}{!}{% Generated by paper/tables.py from results/m13_tables.md; do not edit.
\begin{tabular}{llrrl}
\toprule
\textbf{component} & \textbf{metric} & \textbf{with} & \textbf{without} & \textbf{note} \\
\midrule
support (in WFC propagation) & stable \% (support rule) & 100.0\textsuperscript{\,m6} & 7.8\textsuperscript{\,m6} & off = plain WFC, checked afterwards \\
inventory (in WFC propagation) & valid \% & 100.0\textsuperscript{\,m9} & 1.5\textsuperscript{\,m9} & off = post-filter \\
target shape & cells outside the shape & 0.0\textsuperscript{\,m11} & -- & off: not run \\
balance (no tip-over) & stable \% (physics, table) & 100.0\textsuperscript{\,bl} & 70.2\textsuperscript{\,bl} &  \\
one piece (connect repair) & one piece \% & 100.0\textsuperscript{\,cn} & 3.0\textsuperscript{\,cn} &  \\
one piece, in the demo pipeline & one piece \% & 100.0\textsuperscript{\,dm} & 0.0\textsuperscript{\,dm} &  \\
learned contact style & size-matched judge & $-$0.3281\textsuperscript{\,m10} & $-$0.3721\textsuperscript{\,m10} & off = uniform part/color weights \\
learned proposal & size-matched judge & \textit{pending}\textsuperscript{\,pr} & \textit{pending}\textsuperscript{\,pr} & both raster order, statistics in both \\
raster visit order (control) & size-matched judge & \textit{pending}\textsuperscript{\,pr} & \textit{pending}\textsuperscript{\,pr} & a proposal gain must exceed this \\
\bottomrule
\end{tabular}
}
  \caption{Each constraint or component switched off.}
  \label{tab:ablation}
\end{table*}

\subsection{Constraints in propagation versus post-filtering}
Filtering afterwards does not scale. Plain WFC on plates is stable by the support rule in 19.6\% of builds on a
$4\times4\times2$ grid, 3.9\% at $4\times4\times3$ and in none of 1000 at $8\times8\times6$; with the support
constraint in propagation, 100\% at every size.
% src: results/m6_stability.md (A_plain 19.6 / 3.9 / 0.0; B_aware 100.0)
For a set's inventory, a post-filter keeps 1.5\% of builds (0\% on 9 of 10 sets), while the inventory constraint
with backtracking keeps 100\% and uses 93.1\% of the parts.
% src: results/m9_inventory.md (Findings 1-2)

\subsection{WFC versus the bottom-up decoder}
The bottom-up decoder is valid by construction as well, and is better on every measure we took.
On the support rule it is 100\% stable and 27$\times$ faster than constrained WFC at $8\times8\times6$ (0.020\,s vs
0.533\,s).
% src: results/m6_stability.md (Finding 4; D 0.020 s, B 0.533 s)
On a table, 93.2\% of its builds can stand with some clutch versus 85.6\% for WFC, and 27.1\% need no clutch at
all versus 2.9\%: WFC's overhangs make its builds depend on grip. Every build that cannot stand tips over; on a
baseplate all of them hold.
% src: results/m7_physics.md (8x8x6: any clutch D 93.2 / B 85.6; no clutch D 27.1 / B 2.9; Finding 2)
With a set's plates as inventory it uses 99.9\% of the parts versus 93.1\%, 99.0\% stand versus 83.0\%, at
0.014\,s versus 1.41\,s per build.
% src: results/m9_inventory.md (Finding 3)
With the full catalog, D and D2 build every one of 10 sets in all 20 seeds (up to 746 parts in 4.4\,s), using
99.5\% of the parts; constrained WFC finishes 4 of 20 builds within 30\,s on one set, none on seven, and cannot run
on two for memory. Two causes are visible in its counters: late dead ends from the global part counts (about
1000 backtracks per build) and the raw cost of propagation over large tile sets.
% src: results/m9b_inventory_catalog.md (Findings 2, 4; per-set table: 10015-1 4/20, 10019-1 / 10134-1 not run)
\todo{budget curve: builds finished within 30/120/300/600 s (\texttt{experiments/budget\_curve.py}), Fig.~\ref{fig:budget}}

\paragraph{Target shapes.} With a shape and an inventory together, all 3000 builds of every generator are valid.
WFC fills more of the shape than D (73.0\% vs 68.0\%) because it considers every placement covering a cell, but D2,
which does the same greedily, fills 77.7\%, uses more of the parts and is 20$\times$ faster (0.031\,s vs 0.62\,s).
WFC keeps an edge only on large overhangs (mushroom: 53.2\% vs 49.9\%), where its extra fill relies on long
cantilevers held by grip.
% src: results/m11_shape_inventory.md (mean table: B 73.0, D 68.0, D2 77.7; time 0.619 / 0.031; mushroom 53.2 / 49.9)
% src: results/m11_shape_quick.md (Finding 3: WFC overhang fill relies on grip)

\subsection{Learned style}
With 4-fold cross-validation over the 28 sets (statistics from the 21 training sets), learned contacts bring D2's
builds from a set's own parts closer to that held-out set's statistics on 23 of 28 sets (divergence 0.330 to 0.298),
and shrink the size-matched gap to the real model from 0.073 to 0.029; for WFC, 0.467 to 0.373 on 12 of 13 sets.
Under the judge, D2 beats WFC on 13 of 13 sets at about 1/100 of the time.
% src: results/m10_style.md (Findings 1, 3; tables "All conditions" and "WFC vs bottom-up")

\subsection{The full pipeline}
Two global failure modes appeared once single parts were sound. First, tip-overs: on 10182-1, 186 of 200 candidates
needed infinite clutch, all from small separate stacks; the balance mask takes the physics pass rate from 7\% to
100\% (6372-1: 76\% to 100\%) without changing utilization or the judge. Second, piles: builds fell apart into many
clusters when lifted; the repair pass makes 100\% of builds one piece (from 0.5\% and 5.5\%).
% src: results/m12_pipeline.md (step 1b diagnosis 186/200; m12_balance_summary.json slab 7.0 -> 100, 6372-1 76 -> 100)
% src: results/m12_connect_summary.json (balance 0.5 / 5.5 -> balance_connect 100)
On a spaceship target built from six sets (1978--2010, 203--1514 parts), every D2 candidate is valid, stands, and is
one piece, with 52--74\% fill and 66--92\% utilization; the one-piece rule costs 2--20 points of fill, most on
inventories of 1-wide bricks. WFC finishes none of 50 candidates within the cap (Fig.~\ref{fig:demo}).
% src: results/m12_demo_summary.json (d2 rows: fill 52.3-74.3, utilization 66.1-91.6, valid/physics/connected 100; parts 203-1514)
% src: results/m12_pipeline.md (step 3: "costs 2-20 points of fill", "0 of 50 candidates")
Among the verified candidates of a run, 71\% are distinct designs on average, measured as the effective number of
distinct designs (similarity-sensitive diversity of order 1 \cite{leinster2012diversity}) where two designs are
distinct when at least half a typical D2 distance apart in feature space; the diversity-aware selection keeps the top-$k$ fully distinct.
% src: results/m13_diversity.md (pipeline: effective share 71.4%, selected 100%)

\begin{figure*}[t]
  \centering
  \figorplaceholder{../results/m12_demo_sheet.png}{0.95\textwidth}{six-set spaceship demo sheet}
  \caption{One target (a spaceship silhouette with height and color maps), built from the parts of six official sets
    (1978--2010). Each build is valid, stands on a table and is one piece.}
  \label{fig:demo}
\end{figure*}

\begin{figure}[t]
  \centering
  \figorplaceholder{../results/m12b_before_after.png}{\columnwidth}{3D target before/after (M12b step 1a: heightfield vs exact cells)}
  \caption{A BrickGPT structure as target: solid columns (before) and exact cells (after). \todo{numbers from the 1a write-up}}
  \label{fig:beforeafter}
\end{figure}

\begin{figure}[t]
  \centering
  \figorplaceholder{../results/m12b_catalog_before_after.png}{\columnwidth}{catalog before/after (slopes and corners, M12b-4)}
  \caption{The same build with the rectangular catalog and with slopes and corner parts. \todo{pending M12b-4}}
  \label{fig:catalog}
\end{figure}

\begin{figure}[t]
  \centering
  \figorplaceholder{../results/m13_budget_curve.png}{\columnwidth}{budget curve (experiments/budget\_curve.py)}
  \caption{Share of builds finished versus time budget: constrained WFC and D2 on the same inventories.
    \todo{run on an idle machine}}
  \label{fig:budget}
\end{figure}

\todo{user study: pairwise preferences, Bradley--Terry \cite{bradley1952rank} (experiments/user\_study.py)}

% ------------------------------------------------------------------------------------------------------------
\section{Limitations}
\begin{itemize}
  \item \textbf{Local statistics.} The style proposal and judge see contacts between neighbouring parts only; global
    structure (symmetry, rooms, facades) is not modelled. The learned proposal is meant to address this.
  \item \textbf{Targets.} Model and image targets are heightfields by default (solid columns); exact cells keep
    overhangs and voids but are filled less, and some exact targets produce candidates of clearly different heights.
    % src: results/m12_pipeline.md (step 2: columns 76-84% vs cells 37-49% fill); results/m13_diversity.md (bimodal distances on two BrickGPT-target setups)
  \item \textbf{Support outside the object.} The decoder needs a stud below every part, so overhangs need support
    columns, which a target shape may not contain; the physics check does not reward clever cantilevers.
  \item \textbf{Catalog gaps.} Slopes, corner parts, jumpers and Technic are not yet parts; with the rectangular
    catalog the builds use 35--67\% of each set's parts.
    % src: results/m9b_inventory_catalog.md (Finding 1)
    \todo{update after M12b-4}
  \item \textbf{Physics model.} Vertical forces only, rigid parts, an uncalibrated clutch limit.
  \item \textbf{WFC memory and cost.} Constrained WFC's backtracking history grows with cells $\times$ tiles (e.g.
    2.4\,GB for a 746-part set), so it does not run on the largest inventories at all.
    % src: results/m9b_inventory_catalog.md (10019-1: "not run (memory: 2442 MB)")
\end{itemize}

% ------------------------------------------------------------------------------------------------------------
\section{Conclusion}
Alternate builds from a real inventory can be generated with support, inventory, shape, balance and connectivity all
guaranteed by construction, from any set number, on a CPU in seconds to about two minutes per candidate
depending on the inventory's size.
% src: results/m12_demo_summary.json (seconds_per_candidate 10.6-139.9, 10 in parallel); results/m9b_inventory_catalog.md (D2 4.4 s for 746 parts, single process) Pushing the constraints into
WFC's propagation works, but a bottom-up decoder with the same constraints as masks is simpler, faster and more
robust, so we use WFC as a compared method. Style learned from official sets moves the builds measurably toward real
designs. \todo{learned proposal and user study}

\paragraph{Acknowledgements.} The generate--filter--judge--select framing was inspired by an essay at
\url{https://content.moldandyeast.com/forage/y6fmx6n4}. LDraw models are from the LDraw OMR; inventories from
Rebrickable; BrickGPT's structures from StableText2Brick (MIT).

{\small
\bibliographystyle{plain}
\bibliography{refs}
}

\end{document}
